% 自写模板 — 全国大学生数学建模竞赛 (CUMCM) 论文骨架
% 用法: xelatex main.tex (连编两遍使引用稳定); 或
%       python scripts/render_paper.py --competition cumcm --workspace <paper_workspace>
% 依据: 竞赛公开发布的论文格式规范 (A4、页边距不小于 2.5cm、首页摘要、
%       摘要独占一页、正文自第 2 页起、页码自摘要页起算)。
%       本文件为独立实现, 未参考任何第三方模板源码。
% 字体: 正文宋体小四, 章节标题黑体; xelatex 下 ctex 自动选用系统字体
%       (Windows: 宋体/黑体; Linux 需自备字体或改用 fontset=fandol)。

\documentclass[zihao=-4,a4paper]{ctexart}

% ---------- 页面: A4, 四边距不小于 2.5cm ----------
\usepackage{geometry}
\geometry{a4paper, top=2.5cm, bottom=2.5cm, left=2.5cm, right=2.5cm}

% ---------- 数学与公式编号 (1)(2) ----------
\usepackage{amsmath,amssymb}
% ctexart 默认公式编号即为 (1)(2) 样式, 无需 \numberwithin

% ---------- 图表 ----------
\usepackage{graphicx}
\usepackage{float}      % [H] 就地放置
\usepackage{booktabs}   % 三线表 \toprule \midrule \bottomrule
\usepackage{multirow}
\usepackage{longtable}
\usepackage{caption}
\usepackage{subcaption}

% ---------- 代码附录 ----------
\usepackage{listings}
\usepackage{xcolor}

% ---------- 其他 ----------
\usepackage{enumitem}
\usepackage[hidelinks]{hyperref}

% 章节标题: 黑体 (ctexart 默认即黑体, 此处显式声明字号便于调整)
\ctexset{
  section = {format = {\Large\bfseries\heiti}},
  subsection = {format = {\large\bfseries\heiti}},
  subsubsection = {format = {\bfseries\heiti}},
}

% 图表标题: 黑体小五号 "图1 / 表1"
\captionsetup{font={small,bf}, labelsep=quad}

% 代码样式 (示例代码保持 ASCII, 避免 listings 中文兼容问题;
% 如需在代码中放中文注释, 建议改用 minted 或伪代码环境)
\lstset{
  basicstyle=\small\ttfamily,
  breaklines=true,
  frame=single,
  numbers=left,
  numberstyle=\tiny\color{gray},
  keywordstyle=\color{blue},
  commentstyle=\color{green!50!black},
  stringstyle=\color{red!60!black},
  showstringspaces=false,
  columns=flexible,
}

\begin{document}

% 页码自摘要页起算 (摘要页 = 第 1 页)
\pagenumbering{arabic}

% ============================================================
% 第 1 页: 摘要 (独占一页, 不与正文同页)
% ============================================================

\begin{center}
  {\Large\heiti 论文标题: [论文标题]}\\[0.3em]
  {\large 题号: [题号]\qquad 参赛队号: [参赛队号]}
\end{center}

\vspace{0.5em}

% 摘要标题: 黑体居中
\begin{center}
  {\Large\heiti 摘\quad 要}
\end{center}

% ---- 摘要正文 (公开规范建议 1 页以内, 突出方法/结果/结论) ----
针对 [问题背景一句话重述], 本文建立了 [模型名] 并给出求解方案。
首先, 针对 [子问题 1], 采用 [方法] 建立 [模型], 求解得到 [关键结果];
其次, 针对 [子问题 2], 考虑 [约束/目标], 构建 [模型], 通过 [算法] 求得
[关键结果]; 最后, 对模型进行了灵敏度分析与稳健性检验, 结果表明
[结论一句话]。

本文的特色在于 [建模思路或求解方法的亮点, 一到两句]。

\vspace{1em}
\noindent{\heiti 关键词:} [关键词1]; [关键词2]; [关键词3]; [关键词4]

\newpage

% ============================================================
% 正文 (自第 2 页开始)
% ============================================================

\section{问题重述}

\subsection{问题背景}

[用两到三段交代问题背景与研究意义, 不照抄题面。]

\subsection{问题提出}

\begin{itemize}
  \item [子问题 1 的任务描述];
  \item [子问题 2 的任务描述];
  \item [子问题 3 的任务描述]。
\end{itemize}

\section{问题分析}

[分析每个子问题的本质、变量关系与求解思路, 说明技术路线。]

\section{模型假设}

\begin{enumerate}
  \item [假设 1 及其合理性依据];
  \item [假设 2 及其合理性依据]。
\end{enumerate}

\section{符号说明}

% 三线表示例: 符号说明表 (booktabs)
\begin{table}[H]
  \centering
  \caption{符号说明表}
  \begin{tabular}{cll}
    \toprule
    符号 & 含义 & 单位 \\
    \midrule
    $i$   & 决策变量编号 & --- \\
    $x_i$ & 第 $i$ 个决策变量 & [单位] \\
    $\lambda$ & 模型参数 & [单位] \\
    \bottomrule
  \end{tabular}
\end{table}

\section{模型建立与求解}

\subsection{[子问题 1] 的模型建立}

[模型推导过程。示例公式如下, 编号自动为 (1)(2)。]

决策目标为最小化总成本:
\begin{equation}
  \min \; Z = \sum_{i=1}^{n} c_i x_i,
\end{equation}
约束条件为
\begin{equation}
  \begin{cases}
    \displaystyle\sum_{i=1}^{n} a_{ij} x_i \le b_j, & j = 1,2,\dots,m, \\
    x_i \ge 0, & i = 1,2,\dots,n.
  \end{cases}
\end{equation}

\subsection{模型求解}

[求解算法与步骤说明。]

% 三线表示例: 结果表
\begin{table}[H]
  \centering
  \caption{[求解结果汇总]}
  \begin{tabular}{cccc}
    \toprule
    方案 & 指标 A & 指标 B & 目标值 \\
    \midrule
    方案 1 & [值] & [值] & [值] \\
    方案 2 & [值] & [值] & [值] \\
    方案 3 & [值] & [值] & [值] \\
    \bottomrule
  \end{tabular}
\end{table}

% 图占位示例 (编译无需真实图片; 实际使用时取消注释并替换路径)
\begin{figure}[H]
  \centering
  % \includegraphics[width=0.8\textwidth]{figures/result.png}
  \fbox{\parbox[c][6cm][c]{0.78\textwidth}{\centering [图占位: 技术路线图 / 结果可视化]}}
  \caption{[图题: 一句话说明该图支撑的结论]}
\end{figure}

\subsection{[子问题 2] 的模型建立与求解}

[按同样结构展开。]

\section{灵敏度与稳健性分析}

[对关键参数做扰动实验, 用表格或图给出结果, 并给出结论。]

\section{模型评价与推广}

\subsection{模型的优点}

\begin{itemize}
  \item [优点 1];
  \item [优点 2]。
\end{itemize}

\subsection{模型的缺点}

\begin{itemize}
  \item [缺点 1]。
\end{itemize}

\subsection{模型的改进与推广}

[改进方向与应用场景。]

% ============================================================
% 参考文献 (按 GB/T 7714 顺序编码制手工排列)
% ============================================================
\begin{thebibliography}{9}
  \bibitem{ref1} 作者. 论文题目[J]. 期刊名, 2024, 20(3): 1-10.
  \bibitem{ref2} 作者. 书名[M]. 北京: 出版社, 2023.
  \bibitem{ref3} Author A. Paper title[J]. Journal Name, 2022, 15(2): 100-110.
\end{thebibliography}

% ============================================================
% 附录: 程序代码 (listings; 放不进正文的关键代码)
% ============================================================
\appendix

\section{程序代码}

% 示例: Python 求解脚本骨架 (代码保持 ASCII, 中文注释请在源码中酌情处理)
\begin{lstlisting}[language=Python, caption={求解脚本示例 solve.py}]
import numpy as np
from scipy.optimize import linprog

# model parameters
c = np.array([2.0, 3.0, 1.5])          # cost coefficients
A_ub = np.array([[1.0, 1.0, 0.0],
                 [0.0, 2.0, 1.0]])
b_ub = np.array([40.0, 50.0])

res = linprog(c, A_ub=A_ub, b_ub=b_ub,
              bounds=[(0, None)] * 3, method="highs")
print("optimal value:", res.fun)
print("solution:", res.x)
\end{lstlisting}

\end{document}
